\documentclass[11pt]{article}
\usepackage[margin=1in]{geometry}
\usepackage{booktabs}
\usepackage{array}
\usepackage{hyperref}
\usepackage[T1]{fontenc}
\title{11-768 HW2: Evaluating a Data-Visualization Agent}
\author{Shengzhong Guan}
\date{}

\begin{document}
\maketitle

\section{Part 1: Examine Seed Agent Runs}

\subsection{Comparison of Ground-Truth and Validator Labels}
The baseline performed poorly overall (macro-MCC 0.0268). Table~\ref{tab:seed}
shows the family scores. Its execution judge produced 33 false positives and
eight false negatives. It predicted no \texttt{wrong\_data} or
\texttt{hard\_to\_read} errors, and found only three of 71 evaluated
\texttt{wrong\_chart} positives.

\begin{table}[h]
\centering
\caption{Seed-set MCC from \texttt{output/baseline\_predictions.json}, \texttt{pre\_iteration\_seed\_predictions.json}, and \texttt{final\_seed\_predictions.json}.}
\label{tab:seed}
\begin{tabular}{lrrrrr}
\toprule
Validator & Execution & Wrong data & Wrong chart & Hard to read & Macro \\
\midrule
Baseline & -0.0022 & 0.0000 & 0.1096 & 0.0000 & 0.0268 \\
Part 2 & 1.0000 & 0.4251 & 0.3088 & 0.0301 & 0.4410 \\
Final & 1.0000 & 0.3906 & 0.2975 & 0.0301 & 0.4296 \\
\bottomrule
\end{tabular}
\end{table}

\subsection{Recurring Error Type 1: Contradictory Execution Decisions}
The baseline asked the VLM whether execution failed and parsed the first
\texttt{YES}/\texttt{NO} token. In runs 004, 022, and 023, its own evidence said
that the agent completed and produced \texttt{figure.png}, yet it emitted
\texttt{execution\_failure}. Because this family caused an early return, the
false positive also suppressed the other three judges; run 004, for example,
was human-labeled with both wrong data and wrong chart. Conversely, all twelve
true execution failures (001, 010, 053, 065, 089, 096, 103, 145, 168, 171,
181, and 208) lacked the required root \texttt{figure.png}, but the baseline
missed eight of them.

\subsection{Recurring Error Type 2: Broad Prompts Miss Visual Defects}
The baseline's short independent questions did not force a requirement-by-
requirement comparison. Run 008 contained nearly white node labels on a white
background and omitted a required orange source label. Run 009 left dense date
ticks horizontal rather than at the requested 30-degree rotation, producing an
overlapping band. Run 025 lacked the requested three-level, region-colored
hierarchy and had extensive label collisions. Run 030 rendered one requested
series white on white and replaced the exact suptitle. The baseline accepted
all four.

\subsection{Proposed Changes}
First, replace model-based execution classification with deterministic checks
for the required file, PNG format, and successful decoding. Second, make one
joint structured VLM decision for the remaining families using the complete
rubric, explicit family boundaries, the task, input contents, trajectory, and
rendered image. Require concrete evidence for every positive decision.

\section{Part 2: Improve the Validator}

\subsection{Implemented Changes}
The implementation follows both proposals. \texttt{judge\_execution} uses
Pillow to reject a missing, corrupt, or non-PNG output and otherwise returns an
empty list, allowing later checks to run. The nonterminal judge sends a bounded
excerpt of input and trajectory evidence plus a downscaled image and original
pixel/DPI metadata in one request. Its prompt enumerates wrong values,
selection, ordering, chart constraints, and concrete readability failures. It
expects a JSON object containing one Boolean and evidence string per family,
and validates that schema before constructing errors. This also avoids the
baseline's brittle interpretation of a leading word that contradicts the
explanation.

\subsection{Seed-Set Performance and Qualitative Analysis}
The full Part 2 results appear in Table~\ref{tab:seed}. Execution reached
TP=12, TN=98, FP=0, and FN=0, eliminating both contradictory decisions and the
terminal cascade. Wrong-data MCC rose from 0 to 0.4251 and wrong-chart MCC from
0.1096 to 0.3088. The overall macro-MCC rose from 0.0268 to 0.4410. However,
readability remained the main weakness: the judge found only one of 39
positives (MCC 0.0301).

The revised rubric correctly recovered multiple defects in runs 004, 025, and
030. A representative false negative is run 008: the image visibly contains
near-white labels on white, but the model returned only \texttt{wrong\_chart}
and missed \texttt{hard\_to\_read}. A representative false positive is run
013, where it called requested axis-limit/design discrepancies
\texttt{wrong\_data}; this suggests that the family boundary needs more
emphasis. Run 031 is another false positive: the model's own explanation first
claims the legend mapping is wrong and then says the same mapping is correct,
showing that structured output alone does not eliminate visual hallucination.

\section{Part 3: Design New Evaluation Tasks}

\subsection{Anticipated Generalization Weaknesses}
I targeted two classes. \emph{Derived-data encoding} tests whether the
validator can verify calculations that are not directly readable from the
prompt image alone, especially when several visual channels encode different
derived quantities. \emph{Precise chart constraints} tests whether it can
audit exact scales, styles, subplot relationships, annotations, and layout
instead of merely recognizing the broad chart type.

\subsection{Task Design}
\begin{itemize}
  \item \texttt{regional\_profit\_filter}: filter records, compute and aggregate
  profit, sort descending, and draw labeled horizontal bars.
  \item \texttt{weighted\_bubble\_metrics}: encode latency, accuracy, request
  count, and cost simultaneously through position, area, and color.
  \item \texttt{monthly\_conversion\_combo}: derive conversion rates and a
  cumulative series in stacked panels, including a twin axis and annotations.
  \item \texttt{log\_errorbar\_threshold}: combine log scales, proportional
  error bars, exact limits, shading, grids, and an arrow annotation.
  \item \texttt{shared\_panel\_functions}: enforce four exact functions in a
  shared-axis 2-by-2 layout with panel-specific styles and spine rules.
\end{itemize}
Every task's inputs are empty: all data are synthetic, written by me, and included directly in the prompts, making the
expected values independently checkable without external data or packages.
Each task was run once with Qwen2.5-Coder-3B, Ministral-3-14B, and
GLM-4.7-Flash through the fixed scaffold.

\subsection{Agent Runs, Human Labels, and Validator Performance}
I reviewed each task instruction, complete trajectory, and final figure,
assigned an empty list only when it met the rubric, and treated a
missing/invalid figure as terminal. I also read each PNG's pixel dimensions.
This caught an easy-to-miss failure: several agents created the requested
Matplotlib canvas but saved with \texttt{bbox\_inches="tight"}, changing the
final required pixel dimensions. Those runs are labeled \texttt{wrong\_chart},
consistent with the rubric's exact figure-size requirement.

The Part 2 validator retained perfect execution classification but had
wrong-data MCC -0.2000 because it missed both horizontal-bar ordering errors
and hallucinated two other data errors. It detected eight of eleven chart
errors without a false positive (MCC 0.4264). Readability MCC was -0.0909: it
missed the real overlapping twin-axis labels in
\texttt{monthly\_conversion\_combo\_\_qwen} and hallucinated clipping in a
readable grid.

The agent failures and validator failures were distinct. Qwen's three missing
figures were agent execution failures; its incomplete monthly annotations and
the two non-Qwen \texttt{barh} reversals were also genuine agent mistakes.
Tight-bounding-box saves produced genuine size violations across otherwise
strong figures. By contrast, claims that correct labels were missing or that a
correct cubic series used wrong values were validator hallucinations. The
former require better generation behavior; the latter motivate conservative,
grounded validator checks.

\begin{table}[h]
\centering
\caption{Authored-set MCC.}
\label{tab:authored}
\begin{tabular}{lrrrrr}
\toprule
Validator & Execution & Wrong data & Wrong chart & Hard to read & Macro \\
\midrule
Part 2 & 1.0000 & -0.2000 & 0.4264 & -0.0909 & 0.2839 \\
Final & 1.0000 & 0.6325 & 1.0000 & 1.0000 & 0.9081 \\
\bottomrule
\end{tabular}
\end{table}

\paragraph{Ambiguous labels.} Two judgments were close calls. A figure saved with
\texttt{bbox\_inches="tight"} meets the Matplotlib canvas request but not the final pixel size; I labeled it
\texttt{wrong\_chart} because the task states the saved dimensions. Overlapping twin-axis tick labels in
\texttt{monthly\_conversion\_combo\_\_qwen} were labeled \texttt{hard\_to\_read} because the text is present but difficult to read.

\subsection{Harbor Packaging}
Each task asks for \texttt{plot.py}, \texttt{figure.png}, and a task-specific
\texttt{plotted\_values.json}. Its deterministic tests check file presence,
PNG dimensions, literal expected sidecar values, and reproducibility after
rerunning \texttt{plot.py}. All five reference solutions scored 1.0.
The \texttt{wrong\_total} mutant changes West's computed profit from 90 to 88;
it scored 0 because \texttt{test\_plotted\_values\_are\_exact} rejected the changed value. The
\texttt{illegible} mutant preserves the values and dimensions but uses 2-point
text; it scored 1 because the verifier does not measure rendered legibility.

For these tasks, execution failure and many exact data values can be decided
deterministically. Some requested chart properties can also be checked when
represented in code or metadata, but arbitrary visual design errors and
\texttt{hard\_to\_read} failures generally require image interpretation.
Moreover, a per-task verifier contains hand-derived answers and therefore
cannot grade private runs from unseen tasks; the submitted validator must infer
requirements from each new prompt and its evidence.

\section{Part 4: Iterate on the Design}

\subsection{Additional Improvement Motivated by Part 3}
Part 3 exposed two requirements that were more reliable to verify from metadata
and code than from pixels. First, the joint VLM missed exact saved-image sizes.
I added a conservative parser for explicit ``width by height inches at DPI''
requirements and compare their computed pixel dimensions with the decoded PNG.
Second, both non-Qwen agents sorted values descending before calling
\texttt{barh} but did not reverse the list or call \texttt{invert\_yaxis};
Matplotlib therefore displayed the largest value at the bottom. I added a
trajectory check only when the prompt explicitly requires a descending
horizontal chart from top to bottom and the uncorrected code pattern is present.
Both checks emit grounded evidence and merge with, rather than replace, the
model judgment.

\subsection{Final Evaluation and Qualitative Error Analysis}
On the seed set, the final validator scored macro-MCC 0.4296 (Table~\ref{tab:seed}), slightly below Part 2 (0.4410): the two new checks key on phrases specific to my authored prompts and rarely fire there, so the difference reflects sampling variation in the VLM's answers (17 of 110 runs changed prediction between the two sweeps, with both gains and losses, e.g.\ 191 gained a correct \texttt{hard\_to\_read} while 004 and 050 lost a correct \texttt{wrong\_data}). The seed set therefore shows no regression attributable to the deterministic checks, but also no improvement; the gains below are on the authored set.

On the authored set, macro-MCC increased from 0.2839 to 0.9081. Wrong-data MCC
rose from -0.2000 to 0.6325 because the deterministic \texttt{barh} check found
both previously missed ordering errors. Wrong-chart MCC reached 1.0 because the
dimension check found all three previous false negatives without adding a false
positive. In this final run, the model also correctly classified the one
readability-positive run, so that family reached MCC 1.0; because its evidence
focused on a clipped annotation rather than the most salient overlapping twin
ticks, I do not treat this small-set result as proof of robust readability.

The final authored set has two wrong-data false positives and no false
negatives. For \texttt{shared\_panel\_functions\_\_glm}, the judge claimed the
cubic points do not match $[-8,-1,0,1,8]$, although those exact values are
visible and generated by the trajectory. For
\texttt{monthly\_conversion\_combo\_\_qwen}, it claimed the computed
conversion percentages were wrong even though the executed formula was
correct. Both indicate residual VLM visual hallucination; the deterministic
checks improve requirements with safe programmatic interpretations but do not
attempt to replace arbitrary numerical reasoning.

\subsection{Further Changes and Reproduction}
I made no validator changes beyond the two deterministic checks above. Because the model's sampling is not
deterministic, results can vary between sweeps, as the authored-set rerun shows. Experiments were run with
\texttt{python -m validator.runner --artifacts artifacts|runs --out output/...} followed by
\texttt{python -m workflow score}; the outputs are in \texttt{output/}.

\end{document}
